\documentclass[main.tex]{subfiles}


\begin{document}

%from pagenumber to up
% \phantomsection 
% \hypertarget{toc}{}

 

 

% номер секции вручную
\setcounter{section}{37
}
\addtocounter{section}{-1} 

\section{Магнитный поток}


Пусть имеются два проводящих контура\footnote{Контур~--- это фигура, образованная замкнутой линией в пространстве.} в магнитном поле (рис.\,\ref{fig:p157}).

\def\sub{.417}
\begin{figure}[H]%
\hfill
\begin{subfigure}{\sub\linewidth}
\centering
\includegraphics[]{el/41mflux_pic1.pdf}%
% \caption{При $s_1=s_2$ вторая машина затратила меньше времени.}
% \label{fig:p42sub1}
\end{subfigure}%
\hfill
\begin{subfigure}{\sub\linewidth}
\centering
\includegraphics[]{el/41mflux_pic2.pdf}%
% \caption{При $t_1=t_2$ обе машины прошли одинаковое расстояние.}
% \label{fig:p42sub2}
\end{subfigure}%
\hfill\mbox{}
\caption{Два контура в магнитном поле}
\label{fig:p157}
\end{figure}


Контуры~К$_1$ и~К$_2$ (кольца) расположены в однородном магнитном поле~$B$. Плоскость контура~К$_1$ перпендикулярна линиям поля, а через плоскость контура~К$_2$ линии поля проходят \emph{не} перпендикулярно. Из рис.\,\ref{fig:p157} видно, что кольцо~К$_1$ пронизывается б\'ольшим количеством линий магнитного поля, чем кольцо~К$_2$: в таком случае говорят, что магнитный поток через кольцо~К$_1$ больше, чем через кольцо~К$_2$.

\textbf{Магнитный поток} ($\Phi$ [Вб])~--- это характеристика \emph{количества линий} магнитного поля, пронизывающих контур:
\begin{equation}
    \label{41mflux_e1}
    \Phi=BS\cos\alpha,
\end{equation}
где $B$~--- индукция магнитного поля (в котором находится контур), $S$~--- площадь контура, $\alpha$~--- угол между вектором~$\vec B$ и перпендикуляром к плоскости контура.

На рис.\,\ref{fig:p158} изображен пример к определению магнитного потока.

% контур произвольной формы, находящийся в однородном магнитном поле.

    \begin{figure}[H]
    \centering
\includegraphics[]{el/41mflux_pic3.pdf}%
\caption{Контур произвольной формы в магнитном поле}
\label{fig:p158}
\end{figure}

В магнитном поле с индукцией~$\vec B$ расположен произвольный контур площади~$S$ так, что вектор~$\vec B$ образует угол~$\alpha$ с перпендикуляром~П к плоскости контура. Формула~\eqref{41mflux_e1} дает магнитный поток через этот контур: $\Phi=BS\cos\alpha$.

% Так, в ситуации на рис.\,\ref{fig:p157} с учетом определения магнитные потоки~$\Phi_1$ и~$\Phi_2$ через кольца~К$_1$ и~К$_2$ соответственно связаны так: $\Phi_1>\Phi_2$.


























 
\end{document}